\section{Certified finitary presentations}

\subsection{Deterministic presentations}

A finitary presentation describes only the observations of an object that a
finite proof can consume.

\begin{definition}[Deterministic finitary presentation]
A deterministic finitary presentation \(X\) consists of fixed primitive-
recursive data:
\begin{enumerate}
\item a well-formedness predicate \(\Good_X(x)\);
\item an observation function \(\Obs_X(x,p)\), where \(p\) is a precision or
finite query code;
\item a restriction function \(\operatorname{res}_X(p,q,a)\) for weakening an
answer at precision \(p\) to a weaker precision \(q\);
\item a primitive-recursive relation \(q\preceq_X p\);
\item object-language proofs of identity and composition of restriction and of
observation coherence.
\end{enumerate}
\end{definition}

The precision order need not be ordinary numerical order.  It may be a tuple
of a radius, derivative order, denominator cap, finite open cover, degree cap,
or a finite list of requested observables.  All that matters is that its
comparison and joins are primitive recursive.

\begin{example}[Regular real name]
A regular rational name has observations \(q_p\in\Q\) satisfying
\(|q_p-q_{p+1}|\le 2^{-p}\).  Restriction simply discards excess precision.
The presentation need not contain a completed real number; every theorem is a
finite statement about rational approximations and an explicit error budget.
\end{example}

\subsection{Relational presentations}

Deterministic observations are too strong when valid finite witnesses are not
unique.  Root orderings, triangulations, branch choices, local charts, and
finite covers may change discontinuously.

\begin{definition}[Finitary relational presentation]
A relational presentation \(X\) consists of:
\begin{enumerate}
\item a bound \(B_X(x,p)\) on candidate codes;
\item a decidable validity predicate \(\Cand_X(x,p,a)\) for
\(a<B_X(x,p)\);
\item a compatibility relation
\(\Edge_X(x,p,q,a,b)\) between candidates at related precisions;
\item a finite observable \(\Val_X(x,p,a)\);
\item optional diameter and motion bounds controlling variation among valid
or compatible candidates.
\end{enumerate}
\end{definition}

A candidate cloud is therefore a finite relation, not a set-valued object in a
higher-type universe.  Existence means a bounded disjunction.  Compatibility
means a Boolean table.  Backward viability, developed in \cref{sec:relational},
turns these finite tables into coherent finite paths.

\subsection{Contracted transducers}

\begin{definition}[Contracted transducer]
A contracted transducer \(F:X\to Y\) consists of:
\begin{enumerate}
\item a primitive-recursive backward demand
\(\Need_F(x,p_Y)\), specifying the input precision needed for output precision
\(p_Y\);
\item a bounded state transition or finite relation producing output
candidates from valid input candidates;
\item a local verifier and a quantifier-free invariant;
\item an object-language proof that valid input observations at the demanded
precision yield a valid output observation at \(p_Y\);
\item an extensionality proof showing that observationally equivalent inputs
produce equivalent outputs.
\end{enumerate}
\end{definition}

This definition makes precision flow contravariant and data flow covariant.
The reverse demand is part of the theorem, not an informal estimate hidden in
its proof.

\subsection{Composition}

Suppose \(F:X\to Y\) and \(G:Y\to Z\) are deterministic contracted
transducers.  In the simplest, \emph{static-demand} case the demand of \(G\)
depends only on the requested output query (and fixed finite metadata), so
\[
  \Need_{G\circ F}(x,p_Z)
    =\Need_F\bigl(x,\Need_G(p_Z)\bigr).
\]
More generally, \(G\)'s demand may depend on a designated finite summary of its
input.  Then the contract must also state a primitive-recursive preliminary
demand sufficient to compute that summary, followed by a certified upper
envelope for the final demand.  Equivalently, one may join the demands of all
finitely possible guarded branches.  It is not legitimate to write
\(\Need_G(F(x),p_Z)\) if computing that demand already requires an unbounded or
unknown-precision observation of \(F(x)\).  For a finite program DAG,
\cref{sec:dag} packages the static or majorized case into a reverse pass and a
single forward trace.

\begin{theorem}[Sequential contract composition]
The composite of two contracted transducers is a contracted transducer.  Its
local soundness derivation is obtained by substitution and simultaneous cut
from the local derivations of the factors.
\end{theorem}

\begin{proof}
Fix an output query \(p_Z\).  The contract of \(G\) produces a finite demand
on \(Y\); the contract of \(F\) produces the corresponding demand on \(X\).
Run \(F\) and then \(G\).  The first local derivation concludes the hypotheses
required by the second.  In the repository calculus, instantiate both closed
local proofs with \code{Prf.subst}, use \code{Prf.trans} to replace the
intermediate assumptions, and discharge the input contract by implication
introduction.  Since all formula builders are fixed and quantifier free, the
composite verifier remains quantifier free.
\end{proof}

\subsection{Products, sums, and finite families}

Products combine demands by componentwise join.  A finite family of consumers
of the same producer combines demands by maximum or by a primitive-recursive
join operation in the query preorder.  Tagged sums are safe when the tag is
part of the finite input.  When a tag is selected by approximate numerical
data, one must use a three-way guard or a relational output rather than exact
binary branching.

\subsection{Proof-carrying presentations}

A presentation can carry two distinct kinds of proof.

\begin{enumerate}
\item A \emph{certificate verifier} is a fixed Boolean program.  It checks a
finite witness such as a factorization identity, a matrix inverse, a cover, or
a homotopy.
\item A \emph{proof-producing compiler} is a Lean function whose input fields
are local \(\Prf\,\PRAAx\) derivations and whose output is a global derivation.
\end{enumerate}

The first kind makes domain mathematics modular: a deep theorem may be
isolated as a producer of a finite certificate.  The second kind prevents the
framework itself from becoming merely semantic.

\subsection{Adequacy and honesty}

A presentation is adequate for a theorem only if its observations determine
the theorem's conclusion at the claimed precision.  This requirement rules
out several common shortcuts:
\begin{itemize}
\item a metastable window is not a Cauchy name without a tail modulus;
\item a finite cardinality bound is not an enumerating cutoff;
\item a nonzero proposition is not a numerical inverse guard;
\item a multiset of roots is not an ordered list at a collision;
\item a Selmer class is not a rational point on its covering curve.
\end{itemize}

The framework does not insist on stronger data than the theorem needs.  It
instead insists that every consumed datum be named and verified.  This is the
main sense in which the interface is both constructive and reusable.
